% =======================================================================================
% Appendice B: tabelle complete dei risultati
% =======================================================================================
\chapter{Tabelle complete dei risultati}\label{app:results}

\noindent Questa appendice raccoglie le tabelle complete delle campagne, che il
\autoref{chap:results} riporta in forma aggregata. Le tabelle sono \emph{generate} dai file di
risultato delle campagne dalla procedura \texttt{pruning/make\_appendix\_tables.py}, e i totali che
esse riportano coincidono con i valori discussi nel corpo della tesi: la generazione è quindi
anche una verifica di coerenza fra i numeri citati e i dati grezzi, con la stessa logica della
procedura di verifica automatica descritta nell'appendice sulla riproducibilità.

% =======================================================================================
\section{Campagna principale su Bridges}\label{sec:app-bridges}

La \autoref{tab:app-bridges} riporta le \num{50} configurazioni della campagna principale, con la
medesima ricetta di accoppiamento impiegata nei test: per ogni configurazione la baseline e la
variante \str{A5b}, eseguite sullo stesso dataset, alla stessa soglia, con lo stesso seme. I
totali in fondo alla tabella --- \num{69} imputazioni corrette in più e \num{72} errori in meno
--- sono quelli discussi nella \autoref{sec:res-a5b}.

% FILE GENERATO da pruning/make_appendix_tables.py — non modificare a mano.
\footnotesize
\begin{longtable}{@{}llrrrrrr@{}}
\caption{Campagna principale su Bridges, configurazione per configurazione. ``RFD'' sono le regole in ingresso e in uscita dalla potatura, ``rid.''\ la riduzione percentuale, $\Delta$~corr.\ e $\Delta$~err.\ le variazioni rispetto alla medesima configurazione non potata.}\label{tab:app-bridges}\\
\toprule
dataset & soglia & run & RFD in & RFD out & rid.\ (\si{\percent}) & $\Delta$ corr. & $\Delta$ err. \\
\midrule
\endfirsthead
\toprule
dataset & soglia & run & RFD in & RFD out & rid.\ (\si{\percent}) & $\Delta$ corr. & $\Delta$ err. \\
\midrule
\endhead
\midrule
\multicolumn{5}{r}{totale} & & \textbf{+69} & \textbf{-72} \\
\bottomrule
\endlastfoot
bridges\_14 & 9 & 1 & 8961 & 1366 & 84,8 & +1 & -1 \\
bridges\_14 & 9 & 2 & 8961 & 1381 & 84,6 & +0 & +0 \\
bridges\_14 & 9 & 3 & 8961 & 1388 & 84,5 & +0 & +0 \\
bridges\_14 & 9 & 4 & 8961 & 1368 & 84,7 & +0 & +0 \\
bridges\_14 & 9 & 5 & 8961 & 1391 & 84,5 & +0 & +0 \\
bridges\_28 & 9 & 1 & 8961 & 1395 & 84,4 & +0 & -1 \\
bridges\_28 & 9 & 2 & 8961 & 1406 & 84,3 & +0 & +0 \\
bridges\_28 & 9 & 3 & 8961 & 1410 & 84,3 & -2 & -1 \\
bridges\_28 & 9 & 4 & 8961 & 1415 & 84,2 & -1 & +0 \\
bridges\_28 & 9 & 5 & 8961 & 1427 & 84,1 & -1 & +0 \\
bridges\_42 & 9 & 1 & 8961 & 1450 & 83,8 & +2 & -1 \\
bridges\_42 & 9 & 2 & 8961 & 1548 & 82,7 & +8 & -5 \\
bridges\_42 & 9 & 3 & 8961 & 1545 & 82,8 & +1 & -2 \\
bridges\_42 & 9 & 4 & 8961 & 1566 & 82,5 & +1 & -2 \\
bridges\_42 & 9 & 5 & 8961 & 1433 & 84,0 & -2 & +0 \\
bridges\_56 & 9 & 1 & 8961 & 1718 & 80,8 & +0 & -2 \\
bridges\_56 & 9 & 2 & 8961 & 1854 & 79,3 & -3 & +1 \\
bridges\_56 & 9 & 3 & 8961 & 1877 & 79,1 & +1 & -1 \\
bridges\_56 & 9 & 4 & 8961 & 1568 & 82,5 & -1 & -2 \\
bridges\_56 & 9 & 5 & 8961 & 1628 & 81,8 & +5 & -5 \\
bridges\_70 & 9 & 1 & 8961 & 2248 & 74,9 & +5 & -5 \\
bridges\_70 & 9 & 2 & 8961 & 2317 & 74,1 & +6 & -5 \\
bridges\_70 & 9 & 3 & 8961 & 2362 & 73,6 & +3 & -3 \\
bridges\_70 & 9 & 4 & 8961 & 2266 & 74,7 & +5 & -3 \\
bridges\_70 & 9 & 5 & 8961 & 2808 & 68,7 & +4 & -2 \\
bridges\_14 & 15 & 1 & 18844 & 1853 & 90,2 & +2 & -1 \\
bridges\_14 & 15 & 2 & 18844 & 1881 & 90,0 & +0 & +1 \\
bridges\_14 & 15 & 3 & 18844 & 1891 & 90,0 & +0 & +0 \\
bridges\_14 & 15 & 4 & 18844 & 1849 & 90,2 & +1 & +1 \\
bridges\_14 & 15 & 5 & 18844 & 1896 & 89,9 & -1 & +0 \\
bridges\_28 & 15 & 1 & 18844 & 1904 & 89,9 & +2 & -2 \\
bridges\_28 & 15 & 2 & 18844 & 1918 & 89,8 & +2 & -1 \\
bridges\_28 & 15 & 3 & 18844 & 1923 & 89,8 & -3 & -1 \\
bridges\_28 & 15 & 4 & 18844 & 1934 & 89,7 & -1 & +0 \\
bridges\_28 & 15 & 5 & 18844 & 1961 & 89,6 & -1 & +0 \\
bridges\_42 & 15 & 1 & 18844 & 1994 & 89,4 & +3 & +0 \\
bridges\_42 & 15 & 2 & 18844 & 2139 & 88,6 & +7 & -3 \\
bridges\_42 & 15 & 3 & 18844 & 2158 & 88,5 & +1 & +0 \\
bridges\_42 & 15 & 4 & 18844 & 2177 & 88,4 & +0 & -2 \\
bridges\_42 & 15 & 5 & 18844 & 1974 & 89,5 & -4 & -1 \\
bridges\_56 & 15 & 1 & 18844 & 2438 & 87,1 & +7 & -4 \\
bridges\_56 & 15 & 2 & 18844 & 2663 & 85,9 & -3 & +1 \\
bridges\_56 & 15 & 3 & 18844 & 2707 & 85,6 & -1 & -1 \\
bridges\_56 & 15 & 4 & 18844 & 2200 & 88,3 & -1 & -2 \\
bridges\_56 & 15 & 5 & 18844 & 2253 & 88,0 & +7 & -4 \\
bridges\_70 & 15 & 1 & 18844 & 3524 & 81,3 & +6 & -4 \\
bridges\_70 & 15 & 2 & 18844 & 3472 & 81,6 & +0 & +0 \\
bridges\_70 & 15 & 3 & 18844 & 3672 & 80,5 & +4 & -4 \\
bridges\_70 & 15 & 4 & 18844 & 3413 & 81,9 & +5 & -2 \\
bridges\_70 & 15 & 5 & 18844 & 4502 & 76,1 & +5 & -3 \\
\end{longtable}

\begin{table}[htbp]
\centering
\caption{Ablazione di \str{A5b} sulle 20 configurazioni comuni ai quattro bracci. ``priorità'' è il solo riordino, ``adattiva'' la sola potatura adattiva; la somma delle due variazioni è +5, la variazione della combinazione è +45.}
\label{tab:app-ablation}
\footnotesize
\begin{tabular}{@{}llrrrr@{}}
\toprule
dataset & soglia & $\Delta$ priorità & $\Delta$ adattiva & somma & $\Delta$ \str{A5b} \\
\midrule
bridges\_14 & 9 & +1 & +0 & +1 & +1 \\
bridges\_14 & 9 & +0 & +0 & +0 & +0 \\
bridges\_28 & 9 & +1 & -1 & +0 & +0 \\
bridges\_28 & 9 & +1 & -1 & +0 & +0 \\
bridges\_42 & 9 & +2 & -1 & +1 & +2 \\
bridges\_42 & 9 & +9 & -1 & +8 & +8 \\
bridges\_56 & 9 & +3 & -4 & -1 & +0 \\
bridges\_56 & 9 & +2 & -4 & -2 & -3 \\
bridges\_70 & 9 & +7 & -2 & +5 & +5 \\
bridges\_70 & 9 & +7 & +0 & +7 & +6 \\
bridges\_14 & 15 & +1 & +1 & +2 & +2 \\
bridges\_14 & 15 & +0 & +0 & +0 & +0 \\
bridges\_28 & 15 & +2 & +1 & +3 & +2 \\
bridges\_28 & 15 & +2 & +0 & +2 & +2 \\
bridges\_42 & 15 & +2 & -1 & +1 & +3 \\
bridges\_42 & 15 & +7 & +0 & +7 & +7 \\
bridges\_56 & 15 & +8 & +1 & +9 & +7 \\
bridges\_56 & 15 & +1 & -5 & -4 & -3 \\
bridges\_70 & 15 & -33 & -2 & -35 & +6 \\
bridges\_70 & 15 & +2 & -1 & +1 & +0 \\
\midrule
\multicolumn{2}{r}{totale} & +25 & -20 & +5 & +45 \\
\bottomrule
\end{tabular}
\end{table}

\scriptsize
\begin{longtable}{@{}llrrrr@{}}
\caption{Taratura conservativa della potatura aggressiva ($\alpha = 0{,}99$) sulle 97 configurazioni con baseline disponibile. Riduzione mediana 0,0\si{\percent}, massima 52,9\si{\percent}; variazioni complessive $\Delta$ corr.\ -1, $\Delta$ err.\ -3.}
\label{tab:app-a99}\\
\toprule
dataset & soglia & run & rid.\ (\si{\percent}) & $\Delta$ corr. & $\Delta$ err. \\
\midrule
\endfirsthead
\toprule
dataset & soglia & run & rid.\ (\si{\percent}) & $\Delta$ corr. & $\Delta$ err. \\
\midrule
\endhead
\bottomrule
\endlastfoot
glass\_107 & 3 & 1 & 0,0 & +0 & +0 \\
glass\_107 & 3 & 2 & 0,0 & +0 & +0 \\
glass\_107 & 3 & 3 & 0,0 & +0 & +0 \\
glass\_107 & 3 & 4 & 0,0 & +0 & +0 \\
glass\_107 & 3 & 5 & 0,0 & +0 & +0 \\
glass\_22 & 3 & 1 & 24,1 & +0 & +0 \\
glass\_22 & 3 & 2 & 15,0 & +0 & +0 \\
glass\_22 & 3 & 3 & 23,5 & +0 & +0 \\
glass\_22 & 3 & 4 & 32,3 & +0 & +0 \\
glass\_22 & 3 & 5 & 41,1 & +0 & +0 \\
glass\_43 & 3 & 1 & 1,1 & +0 & +0 \\
glass\_43 & 3 & 2 & 0,4 & +0 & +0 \\
glass\_43 & 3 & 3 & 0,9 & +0 & +0 \\
glass\_43 & 3 & 4 & 1,3 & +0 & +0 \\
glass\_43 & 3 & 5 & 1,1 & +0 & +0 \\
glass\_64 & 3 & 1 & 0,0 & +0 & +0 \\
glass\_64 & 3 & 2 & 0,3 & +0 & +0 \\
glass\_64 & 3 & 3 & 0,5 & +0 & +0 \\
glass\_64 & 3 & 4 & 0,0 & +0 & +0 \\
glass\_64 & 3 & 5 & 0,3 & +0 & +0 \\
glass\_86 & 3 & 1 & 0,1 & +0 & +0 \\
glass\_86 & 3 & 2 & 0,0 & +0 & +0 \\
glass\_86 & 3 & 3 & 0,0 & +0 & +0 \\
glass\_86 & 3 & 4 & 0,0 & +0 & +0 \\
glass\_86 & 3 & 5 & 0,1 & +0 & +0 \\
restaurant\_104 & 3 & 1 & 0,0 & +0 & +0 \\
restaurant\_104 & 3 & 2 & 4,0 & +0 & +0 \\
restaurant\_104 & 3 & 3 & 0,0 & +0 & +0 \\
restaurant\_206 & 3 & 1 & 0,0 & +0 & +0 \\
restaurant\_206 & 3 & 2 & 0,0 & +0 & +0 \\
restaurant\_206 & 3 & 3 & 0,0 & +0 & +0 \\
restaurant\_259 & 3 & 1 & 0,0 & +0 & +0 \\
restaurant\_259 & 3 & 2 & 0,0 & +0 & +0 \\
restaurant\_52 & 3 & 1 & 16,0 & +0 & +0 \\
restaurant\_52 & 3 & 2 & 24,0 & +0 & +0 \\
restaurant\_52 & 3 & 3 & 16,0 & +0 & +0 \\
restaurant\_104 & 6 & 1 & 5,6 & +0 & +0 \\
restaurant\_104 & 6 & 2 & 9,7 & +0 & +0 \\
restaurant\_104 & 6 & 3 & 0,0 & +0 & +0 \\
restaurant\_206 & 6 & 1 & 0,0 & +0 & -1 \\
restaurant\_206 & 6 & 2 & 0,0 & +0 & +0 \\
restaurant\_206 & 6 & 3 & 0,0 & +0 & +0 \\
restaurant\_259 & 6 & 1 & 0,0 & +0 & +0 \\
restaurant\_259 & 6 & 2 & 0,0 & -1 & -1 \\
restaurant\_52 & 6 & 1 & 37,9 & +0 & +0 \\
restaurant\_52 & 6 & 2 & 45,2 & +0 & +0 \\
restaurant\_52 & 6 & 3 & 45,2 & +0 & +0 \\
glass\_107 & 9 & 1 & 0,0 & +0 & +0 \\
glass\_107 & 9 & 2 & 0,0 & +0 & +0 \\
glass\_107 & 9 & 3 & 0,0 & +0 & +0 \\
glass\_107 & 9 & 4 & 0,0 & +0 & +0 \\
glass\_107 & 9 & 5 & 0,0 & +0 & +0 \\
glass\_22 & 9 & 1 & 31,7 & +0 & +0 \\
glass\_22 & 9 & 2 & 17,4 & +0 & +0 \\
glass\_22 & 9 & 3 & 24,8 & +0 & +0 \\
glass\_22 & 9 & 4 & 38,7 & +0 & +0 \\
glass\_22 & 9 & 5 & 52,9 & +0 & +0 \\
glass\_43 & 9 & 1 & 0,4 & +0 & +0 \\
glass\_43 & 9 & 2 & 0,2 & +0 & +0 \\
glass\_43 & 9 & 3 & 0,5 & +0 & +0 \\
glass\_43 & 9 & 4 & 0,6 & +0 & +0 \\
glass\_43 & 9 & 5 & 0,4 & +0 & +0 \\
glass\_64 & 9 & 1 & 0,0 & +0 & +0 \\
glass\_64 & 9 & 2 & 0,1 & +0 & +0 \\
glass\_64 & 9 & 3 & 0,2 & +0 & +0 \\
glass\_64 & 9 & 4 & 0,0 & +0 & +0 \\
glass\_64 & 9 & 5 & 0,0 & +0 & +0 \\
glass\_86 & 9 & 1 & 0,0 & +0 & +0 \\
glass\_86 & 9 & 2 & 0,0 & +0 & +0 \\
glass\_86 & 9 & 3 & 0,0 & +0 & +0 \\
glass\_86 & 9 & 4 & 0,0 & +0 & +0 \\
glass\_86 & 9 & 5 & 0,0 & +0 & +0 \\
glass\_107 & 15 & 1 & 0,0 & +0 & +0 \\
glass\_107 & 15 & 2 & 0,0 & +0 & +0 \\
glass\_107 & 15 & 3 & 0,0 & +0 & +0 \\
glass\_107 & 15 & 4 & 0,0 & +0 & +0 \\
glass\_107 & 15 & 5 & 0,0 & +0 & +0 \\
glass\_22 & 15 & 1 & 34,6 & +0 & +0 \\
glass\_22 & 15 & 2 & 20,3 & +0 & +0 \\
glass\_22 & 15 & 3 & 23,2 & +0 & +0 \\
glass\_22 & 15 & 4 & 36,0 & +0 & -1 \\
glass\_22 & 15 & 5 & 51,6 & +0 & +0 \\
glass\_43 & 15 & 1 & 0,3 & +0 & +0 \\
glass\_43 & 15 & 2 & 0,2 & +0 & +0 \\
glass\_43 & 15 & 3 & 0,5 & +0 & +0 \\
glass\_43 & 15 & 4 & 0,5 & +0 & +0 \\
glass\_43 & 15 & 5 & 0,2 & +0 & +0 \\
glass\_64 & 15 & 1 & 0,0 & +0 & +0 \\
glass\_64 & 15 & 2 & 0,1 & +0 & +0 \\
glass\_64 & 15 & 3 & 0,1 & +0 & +0 \\
glass\_64 & 15 & 4 & 0,0 & +0 & +0 \\
glass\_64 & 15 & 5 & 0,0 & +0 & +0 \\
glass\_86 & 15 & 1 & 0,0 & +0 & +0 \\
glass\_86 & 15 & 2 & 0,0 & +0 & +0 \\
glass\_86 & 15 & 3 & 0,0 & +0 & +0 \\
glass\_86 & 15 & 4 & 0,0 & +0 & +0 \\
glass\_86 & 15 & 5 & 0,0 & +0 & +0 \\
\end{longtable}


% =======================================================================================
\section{Studio di ablazione}\label{sec:app-ablation}

La \autoref{tab:app-ablation} riporta le venti configurazioni comuni ai quattro bracci
dell'ablazione. La tabella rende verificabile la lettura discussa nella \autoref{sec:ablation}: la
somma delle variazioni dei due componenti isolati è pari a $+5$, mentre la variazione della
combinazione è pari a $+45$, e la differenza ($+40$) è la sinergia che il \autoref{sec:mechanism}
spiega.

% =======================================================================================
\section{Taratura conservativa}\label{sec:app-a99}

La \autoref{tab:app-a99} riporta le \num{97} configurazioni della campagna con il parametro di
aggressività portato al valore conservativo. Vi compare il dato su cui poggia la conclusione della
\autoref{sec:res-boundary}: la riduzione mediana è nulla, la massima è pari al
\SI{52.9}{\percent}, e le variazioni complessive sono di una imputazione corretta in meno e di tre
errori in meno. È questa la ragione per cui la strategia, in quella taratura, viene descritta come
\emph{non dannosa} e non come vantaggiosa: l'effetto sulla qualità non è distinguibile da zero
(\autoref{tab:statistics}), mentre quello sulla riduzione è concentrato in poche configurazioni.
